The increasing interest in sustainable modes of transport and cost effective commuting has motivated the creation of grand ride-sharing concepts in the global context. Wheel Share is a community-based priority ride-sharing system which offers a secure, scalable, and easy experience to customers and the passengers who travel in the same direction. Will using Flutter as the frontend and Supabase (PostgreSQL) as the backend and Google API for location. The system promotes car sharing that saves gas and amounts of traffic growth, traffic congestion and encourage the sustainability of the environment. Ride-Sharing in a period when gas prices are increasing and where the issue of environmental awareness is on the rise.
Ride-Sharing has turned to be a viable and greener way of transport. However, most existing ride sharing services are commercialized and have high commissions, which makes it community-based, affordable and centralized alternative. Through this platform, users can is used in order to post available rides.
Passenger can also search for and send offers to join.
In is also provided on the platform built in controls on financial overviews, rating schemes on critiques, and an emergency option which guaranties safety of the user.
It uses cross platform flexibility of Flutter and Supabase real time backend services, Wheel Share makes data management and communication transparent as well as reliable and efficient performance. It does not only create trust and user-to-user collaboration, but it also
is part of the overall objective of sustainable shared mobility
\cite{mitropoulos2021}.


\section{Background}

Transportation remains a daily challenge for millions of individuals, especially in urban areas where congestion, fuel prices, and time delays create significant inconvenience. The idea of carpooling or ride sharing emerged to address these issues by allowing people traveling in the same direction to share a single vehicle, thereby reducing costs and emissions. However, the local communities, in spite of the introduction of the ride sharing services into the global community. We already do not have easy access to low-cost and safe prospects that can address their needs. These challenges have been such that Wheel Share has been created to handle them. The project integrates current mobile and web technologies to turn drivers and a flawless  experience as well for passengers. The user can create profiles, add rides, search, since it has a user-friendly interface to match in journeys, make and receive offers. Both user type, two types of drivers, passenger and administrator is a role that occurs, which has both role and access privileges to assure ordered functionality and data security.
Moreover, it has the emergency alert system, rating, review and capabilities.The earnings spending dashboards also help individuals save time and money besides offering assurance to the users. Wheel Share will,therefore, be a revolution a change of commuting culture to a more ecological and intermediate kind.


\section{Objective}

It will aim at providing a safe convenient along with reliable car-pooling environmentment by integrating the real time communication, open financial track and security features.The main aims will be the creation of a system that will provide the driver and the passengers with the opportunity to communicate and operate rides effectively and with the aim of developing trust, provide feedbacks in a smooth manner\cite{ashkrof2022}. It also dedicated to the empowerment of the users and providing them with the entire control over their rides and trade, and at the same time enable administrators to trace activities and ensure that the platform possesses integrity.

Lastly, Wheel Share will target combination of convenience, cost efficiency and security, provide-
Development of the alternative to the traditional commuting which will be a win-win to the users and the environment.


\section{Problem Description}

Numerous low-cost modes of transportation are being used as a result of the rising need for affordable, effective travel and transportation.Travel costs, expensive forms of transportation, and lack of reliable car-sharing services are all things that people dislike. The market offers a wide range of solutions.The data is heavy, commercial, or geographically restricted, making it impractical for personal or community use.User comments and market observations can help identify some of these.
Through market observation and user feedback, several core issues were identified:


\begin{enumerate}[label=\roman*.]
    \item Lack of localized, affordable ride sharing solutions tailored to community use.
    \item Absence of a transparent system that allows users to view and compare rides based on city, address, and date.
    \item Inadequate financial tracking and transaction monitoring between drivers and passengers.
    \item Safety concerns due to insufficient identity verification and lack of emergency protocols.
    \item Lack of comprehensive dashboards to analyze travel spending and earnings performance.
    

\end{enumerate}


\section{Project Objectives}
The Wheel Share project has been designed with a clear set of objectives to ensure effective implementation and user satisfaction. The objective is to develop a cross-platform mobile application using Flutter that connects drivers and passengers in an efficient, secure, and transparent manner.

The system intends to implement role-based authentication through Supabase to ensure that each user type driver, passenger, and admin has specific access rights and responsibilities. Drivers will be able to post available rides with relevant details and manage offers from users. Passengers will be able to search for rides by city, address, and date, send offers, and manage accepted rides. Additionally, both parties will have access to dashboards that summarize earnings and expenses respectively, helping them maintain financial clarity.

Another key objective is the incorporation of a rating and review system that fosters accountability and builds trust among users. The inclusion of an emergency feature enhances safety by allowing users to alert the admin in case of unforeseen circumstances\cite{gangadharaiah2023}. The admin panel provides tools to monitor user activity, manage rides, handle emergencies, and remove inappropriate users or content.

Together, these objectives ensure that Wheel Share becomes a holistic and user driven platform that encourages shared mobility while maintaining transparency, safety, and convenience.



\section{Project Scope}

The production of a ride sharing application is part of the Wheel Share combines safe authentication, tools of managing and communication cohesive system. The project will be directed at providing a stable platform that contributes to provision posting of rides, searching and managing and ensuring that it is safe and transparent to all users.
The system has three modules namely driver, passenger and administrator. The driver module provides their users with a chance to post rides, offers, and track it, view, and access their ride history. The passenger module allows the ride searching to take place and offers an offer dispatching, and cost tracking. The administration module provides management functionality, monitoring of processes on the platform and management of emergencies and the efficient work of the operation of the system. The user interface is Flutter and the Database and Authentication is Supabase services, Wheel Share will offer data synchrony, data reliability and data efficiency at real time. The concept of analysing it, in the end, all gives to the ride it has made a ride that is both community based and functional Sharing application that promotes price effectiveness, safety and sustainability.
In conclusion, Wheel Share system is capable of realizing the recent technology and community partnership that will provide the new generation, dependable, and green ride-sharing solution. Helping to establish trustful relations and open transactions, and it enhances the experience of commuting to the entirety and generates the commuting experience and 
tend to smarter and better networked cities.


\section{Solution Application Areas}

The Wheel Share system is designed to serve a wide range of users who seek affordable and reliable commuting options. It provides an intelligent ride sharing platform suitable for daily commuters, students, employees, and individuals seeking intercity travel options. The system can be applied in both urban and suburban contexts, where it enables users to share rides, reduce travel expenses, and optimize vehicle usage.

The application also provides significant benefits for environmental sustainability, as it reduces traffic congestion and lowers carbon emissions through shared mobility. Its role-based dashboards empower different stakeholders drivers, passengers, and administrators by providing access to relevant tools and data insights\cite{soria2021}. For drivers, the platform acts as an earnings management tool, while passengers benefit from transparency in travel spending. Administrators can utilize the platform to monitor safety alerts, manage user activities, and ensure smooth operational flow.


\section{Risk Involved}

\subsection{Technical Challenges}
One significant technical challenges in creating Wheel Share is guaranteeing scalability and performance efficiency. The system needs to handle a high volume of ride data while ensuring real-time synchronization among numerous users. Possible risks include problems with integrating Supabase, performance lags in Flutter interfaces, and difficulties in handling multiple concurrent user sessions. These technical challenges, if not swiftly addressed, could affect system reliability and user satisfaction.Therefore, proactive testing, code refinement, and a strong backend design are essential to reduce technical risks.

\subsection{User Adoption}
Another risk is the user adoption. Ride-sharing in which the communication is a significant aspect is also of importance.The most important secret of success is participation that involves the amount of users. Without proper marketing and awareness, the platform will have low participation of drivers or passengers. Local promotions, strategic outreach programs to control this and referrals as incentives can also help to improve the uptake and continued usage of the Wheel Share application
\cite{compostella2025}.

\subsection{Legal Compliance}
The project must also be able to meet the requirements of laws of transportation and data privacy at regional level.Confidence of the legislation of the check of users, protection of data and ride sharing policies should be introduced to avoid the potential legalization. Moreover, maintaining the field of transparency in user identity check and ride histories will help the platform to achieve control and earn customer trust.

 
\clearpage



